
\section{与拉格朗日子流形关联的谱序列}
\label{sec:6-4}
\liftIndex{I368}

\subsection{构造}
\label{subsec:6-4-1}

至此，谱序列是以纯代数方式构造的。
本节把上一节的结果用于几何情形，
以证明定理~\ref{thm:D}。
在本小节证明其中的 (D.1) 与 (D.2)。

设 $L$ 是辛流形 $M$ 的
相对 spin 拉格朗日子流形，
而 $(C(L;\Lambda_{0,\mathrm{nov}}),\mathfrak m)$
是与之关联的滤过 $A_\infty$ 代数。
取
$\mathbf b=(\boldsymbol b,b)
\in\mathcal M_{\mathrm{weak,def}}(L)$。
由第 3.5 节的构造，
$C(L;\Lambda_{0,\mathrm{nov}})$
是第~\ref{subsec:6-3-1} 小节意义下的 c.f.z.。
而
$(C(L;\Lambda_{0,\mathrm{nov}}),
\mathfrak m_1^{\mathbf b})$
是上链复形。
此外，格罗莫夫紧性定理
保证这个复形满足有隙条件
（见定理~\ref{thm:3-5-11}）。

并且，
$(C(L;\Lambda_{0,\mathrm{nov}}),
\mathfrak m_1^{\mathbf b})$
具有定义~\ref{def:6-3-27} 意义下的弱有限性。
事实上，它与
\[
\bigl(H(L;\Lambda_{0,\mathrm{nov}}),
\mathfrak m_1^{\mathrm{can},\boldsymbol b,b'}\bigr)
\]
上链同伦等价，
其中
$b'\in\mathcal M_{\mathrm{weak}}
\bigl(H(L;\Lambda),
\mathfrak m^{\mathrm{can},\boldsymbol b}\bigr)$。
这是定理~\ref{thm:5-4-2} 的推论。

因此，定理~\ref{thm:6-3-28}
给出定理~\ref{thm:D} 的
(D.1)、(D.2) 所述谱序列。
再利用推论~\ref{cor:5-4-21}，
便可类似地构造定理~\ref{thm:6-1-4}
中满足 (6.1.5.1)、(6.1.5.2) 的谱序列。
